\cotefile{46}
\scanpage{24}

\begin{center}
Définition du morphisme Verschiebung.
\end{center}

Je ne sais pas exactement ce que tu veux que je définisse\deleted{\unsure{z|s}}. Cependant, comme tu affirmes avoir perdu mon tapis de l'année dernière, je vais essayer de te le revendre. Pour cela, je vais supposer que le schéma de base est le spectre d'un corps $k$ de caractéristique $p\handwritten{>}0$, parfait ou non : à toi de voir ce qui se passe dans les autres cas.

Je vais commencer par le cas d'un $k$-groupe commutatif \uline{affin\unsure{\deleted{t}e|\deleted{e}t}} $G$, d'anneau des coordonnées $\mathbb{A}$. Les notations sont les suivantes :

\noindent - $\mathbb{A}^{(p)}$ est la $k$-algèbre déduite de $\mathbb{A}$ par l'extension des scalaires $F_k:k\mathrel{\text{-{}-{}-{}-{}-{}-}\!\handwritten{>}}k$, où $\mathbb{F}_k(x)=x^p$. Un élément de $\mathbb{A}^{(p)}$ s'écrira donc $x\handwritten{\otimes}a$, où $x\handwritten{\in}k$ et $a\handwritten{\in}\mathbb{A}$ ; on notera la règle de calcul $x\handwritten{\otimes}y.a=x.y^p\handwritten{\otimes}a$. Si $k$ est parfait, $\mathbb{A}^{(p)}$ s'identifie à la $k$-algèbre déduite de $\mathbb{A}$ par la restriction des scalaires $x\mathrel{\text{-{}-{}-{}-{}-{}-}\!\handwritten{>}}x^{1\handwritten{/}p}$.

\noindent - $F:\mathbb{A}^{(p)}\mathrel{\text{-{}-{}-{}-{}-{}-}\!\handwritten{>}}\mathbb{A}$ est l'application \deleted{\illegible{1}} $x\handwritten{\otimes}a\mathrel{\text{-{}-{}-{}-{}-{}-}\!\handwritten{>}}x.a^p$.

\noindent - $G^{(p)}$ est le groupe algébrique d'anneau des coordonnées $A^{(p)}$ ; on désignera encore par $F$ le morphisme de $G$ dans $G^{(p)}$ qui est défini par l'application $F:\mathbb{A}^{(p)}\mathrel{\text{-{}-{}-{}-{}-{}-}\!\handwritten{>}}A$.

\noindent - $\handwritten{\otimes}^{p}A$ désigne le produit tensoriel de $p$ exemplaires de $A$ ; $\handwritten{\Delta}^{p}$ est l'application de $\mathbb{A}$ dans $\handwritten{\otimes}^{p}\mathbb{A}$ qui est définie par la diagonale $\handwritten{\Delta}$ dela coalgèbre $\mathbb{A}$ ; $m^p$ est
% archive sheet number: 23
